% !TeX root = ../main.tex
\section{Интеграция полосового фильтра в тракт усилителя и совместное электродинамико-схемотехническое моделирование}

На заключительном этапе проектирования была выполнена задача анализа сквозной амплитудно-частотной характеристики (АЧХ) тракта, состоящего из малошумящего усилителя с согласующими цепями и электродинамической 3D-модели диэлектрического резонатора (рис.~\ref{fig:top_schematic}).

\begin{figure}[H]
    \centering
    \includegraphics[width=0.75\textwidth]{06_sch.png}
    \caption{Схема включения электродинамической подсхемы фильтра на выходе усилительного каскада}
    \label{fig:top_schematic}
\end{figure}

Полноволновая модель фильтра \textit{Filter\_DR} была подключена в качестве иерархического $S$-параметрического блока (\textit{FieldSolver Co-simulation}) в разрыв между выходным согласующим контуром каскада ($C_{51} = 10\text{ пФ}$, $L_{50} = 1{,}62\text{ нГн}$) и стандартным нагрузочным портом $R_H = 50\text{ Ом}$.

\subsection{Электродинамический анализ фильтра на диэлектрическом резонаторе}

Перед выполнением совместного схемотехнического расчета были рассчитаны терминальные параметры рассеяния 3D-модели фильтра (\textit{Filter\_3D}) в широком диапазоне частот от $0{,}5$ до $5{,}0\text{ ГГц}$. Полученные частотные зависимости коэффициентов отражения $\text{dB}(S_{t11})$ и передачи $\text{dB}(S_{t21})$ представлены на рис.~\ref{fig:dr_sparam}.

\begin{figure}[H]
    \centering
    \includegraphics[width=0.95\textwidth]{fig2.png}
    \caption{Частотные зависимости терминальных $S$-параметров 3D-модели фильтра на ДР}
    \label{fig:dr_sparam}
\end{figure}

На полученной характеристике четко прослеживаются резонансные моды колебаний диэлектрического резонатора[cite: 15]. В окрестности целевой частоты $4{,}45\dots 4{,}5\text{ ГГц}$ формируется требуемый полосно-пропускающий пик передачи, в то время как вне полосы пропускания обеспечивается глубокое запирание сигнала (подавление превышает $60\dots 100\text{ дБ}$ на низких частотах).

\subsection{Амплитудно-частотная характеристика усилителя с фильтром}

В схемном редакторе \textit{Nexxim} был выполнен совместный расчет линейных параметров (\textit{Linear Frequency Analysis}) тракта «усилитель + фильтр» в диапазоне частот от $1{,}5$ до $5{,}5\text{ ГГц}$. График результирующего коэффициента передачи тракта $\text{dB}(S(\text{Port2}, \text{Port1}))$ представлен на рис.~\ref{fig:lna_filter_s21}.

\begin{figure}[H]
    \centering
    \includegraphics[width=0.95\textwidth]{fig1.png}
    \caption{Сквозная АЧХ коэффициента передачи $\text{dB}(S(\text{Port2}, \text{Port1}))$ малошумящего усилителя с фильтром}
    \label{fig:lna_filter_s21}
\end{figure}

Анализ итоговой частотной характеристики подтверждает успешную интеграцию каскадов:
\begin{itemize}
    \item В полосе пропускания $4{,}45\dots 4{,}55\text{ ГГц}$ коэффициент передачи достигает величины около $+2\text{ дБ}$, обеспечивая прохождение и усиление полезного радиочастотного сигнала;
    \item За пределами полосы пропускания наблюдается резкий спад АЧХ: при отстройке до $4{,}0\text{ ГГц}$ затухание составляет более $-60\text{ дБ}$, а на частотах ниже $2{,}5\text{ ГГц}$ превышает $-110\text{ дБ}$;
    \item Совместная работа усилителя и диэлектрического резонатора обеспечивает высокую избирательность приемного тракта по зеркальному и внеполосным каналам приема.
\end{itemize}

\newpage
\section*{Заключение}
\addcontentsline{toc}{section}{Заключение}

В ходе выполнения лабораторной работы был спроектирован тракт малошумящего усилителя СВЧ-диапазона, интегрированный с полосовым фильтром на диэлектрическом резонаторе:
\begin{enumerate}
    \item Проведен нелинейный анализ устойчивости каскада во временной области (\textit{Transient Analysis}), подтвердивший отсутствие самовозбуждения и периодическую стабильность колебаний при питании $2{,}0\text{ В}$.
    \item Выполнен электродинамический пересчет 3D-модели цилиндрического диэлектрического резонатора с центральной частоты $3{,}5\text{ ГГц}$ на рабочую частоту $4{,}5\text{ ГГц}$ с коэффициентом масштабирования $k \approx 0{,}778$ ($R = 4{,}05\text{ мм}$, $H = 8{,}1\text{ мм}$).
    \item Рассчитана матрица рассеяния электродинамической модели в Ansys HFSS, подтвердившая формирование селективной полосы пропускания на частоте $4{,}45\dots 4{,}5\text{ ГГц}$.
    \item Осуществлена сквозная схемотехническая интеграция усилителя и 3D-модели фильтра в среде Nexxim, получена результирующая АЧХ системы с эффективным усилением в полосе $4{,}5\text{ ГГц}$ и подавлением внеполосных помех свыше $60\text{ дБ}$.
\end{enumerate}

\newpage
\begin{thebibliography}{9}
\addcontentsline{toc}{section}{Список литературы}
\bibitem{kurushin1} Курушин~А.\,А. Гибридное моделирование СВЧ структур в HFSS ANSYS. --- 1--44~с.
\bibitem{bankov} Банков~С.\,Е., Курушин~А.\,А. HFSS ANSYS Электродинамическое моделирование сложных СВЧ структур, 2006~год.
\end{thebibliography}
